\section{Discussion and Limitations}
\label{sec:limits}
\subsection{Industrial-Informatics Interpretation}
The principal result is that one supervisory path carried four different jobs
while keeping command closure, the submitted verdict, task state, and process
attainment as separate records; relative to the nearest neighbors of
Section~\ref{sec:related}, the differentiator remains production-scoped
authority and evidence on a physical write path. In run
\resRun{} all four scenario predicates
hold, and the governed loop reaches $J/J^{*}$ between \resControllerMin{}
and \resControllerMax{}.

\subsection{Validity and Reproducibility}
The snapshot fixes run identifier, seed, commit, harness fingerprint, and raw
artifacts, so the reported counts can be re-derived rather than trusted. Three
qualifications bound what they support. First, the evidence is simulation-tier:
the simulators are deterministic per seed but are not validated digital twins,
and no physical-device campaign, operator study, Tennessee Eastman Process
(TEP)/Secure Water Treatment (SWaT) replay, or external-framework comparison
is reported. Second, the controller optimum and the
scenario predicates are internal oracles, the ablation has three
deterministic repetitions, and initial state, asynchronous sampling, response
windows, and simulator timing still influence deterministic policies, so
results support mechanism attribution rather than
statistical inference; latencies are single-host observations, not service-level
bounds. Third, several
capabilities of Section~\ref{sec:arch} are implemented but not
exercised---memory retrieval quality, plugin hot-reload,
alarm-chain timing, engine interchange---and the live turns share one provider
and model, their memory call
returning no stored entry, so that layer is exercised rather
than validated; the four process models
broaden structural variety, not transfer.

\subsection{Safety Positioning and Human Oversight}
The objection that a language-model team should not ``take over'' production
applies to a design \system{} does not propose: fast regulation, hardware
interlocks, and safety functions remain in the PLC;
Section~\ref{sec:hitl} bounds what an agent can change; and prompts,
retrieved documents, and memory are untrusted text, not authority
\cite{greshake2023injection,liu2024pinjection}---the proposer is confined to
per-binding scope and journaled,
attributable actions: a bounded blast radius, not a protection guarantee.
Binding scope, admission limits,
readback, and the backstop run with or without human approval, bounding
approval fatigue and automation bias
\cite{bainbridge1983ironies,parasuraman2000model,leesee2004trust,parasuramanmanzey2010complacency};
approval is a revocable per-binding policy; because admission is re-checked at
execution, a stale approval cannot carry an out-of-window value nor outlive a
context change.

\subsection{Deployment Boundaries}
Execution and recovery are non-atomic: bookkeeping may fail after actuation, the
post-write hold guards write frequency not stability, readback agreement
need not imply process recovery, and multi-node compensation may partially
succeed; oscillation, storage failure, clock error,
worst-case recovery time, cross-run time-range
queries (records lack run identifiers), and concurrency beyond the four-line
campaign remain unevaluated.

In production, bindings govern sanctioned tools, not arbitrary privileged-engine
I/O, so least privilege, endpoint security, and IT/OT
separation remain necessary \cite{iec62443-3-3,nist80082}, and independent
hardware interlocks and any required IEC~61511 protection must stay in place: \system{}
is neither a safety instrumented function nor a substitute for plant
protection \cite{iec61511}. The approval wait already fails closed on expiry,
but line-scoped approver authorization is still missing, so production
deployment needs segregation of duties and
alarm-path failure tests \cite{eemua191}.

\section{Conclusion}
\label{sec:conclusion}
\system{} contributes a governed supervisory execution architecture rather
than a new process optimizer: one stable path from agent-team identities
through semantic nodes, production context, governed writes, readback,
observations, verdicts, and compensation, exercised by a \resPhases-phase
suite, agent-team missions, a three-seed loop, and four missions---all
\resChecksTotal{} checks passing and all four mission predicates holding in
the reference run. Physical-plant deployment,
operator studies, and external-framework comparison remain open; until then,
claims of optimization superiority, autonomous control, and safety assurance
stay unsupported.

\section*{Data Availability and Disclosure}
\footnotesize
Artifacts are archived in \path{bench/results/}: main run
\path{bench/results/20260921184204-10lc}, ablation companion
\path{bench/results/20260921181234-e1lite}. Reproduction needs
the simulator source and reset state; the revision artifact will pin
the currently missing simulator-tree fingerprint. AI assistants
supported editing and checking under
author direction; the authors remain responsible for all claims and data.
